\documentclass[authoryear,preprint,12pt]{elsarticle}

\usepackage[
  a4paper,
  left=20mm,
  right=20mm,
  top=20mm,
  bottom=20mm
]{geometry}

\usepackage{amssymb}
\usepackage{amsmath}
\usepackage{algorithm}
\usepackage{algpseudocode}
\usepackage{array}
\usepackage{booktabs}
\usepackage{multirow}
\usepackage{makecell}
\usepackage{enumitem}
\usepackage{siunitx}
\usepackage{adjustbox}
\usepackage[table]{xcolor}
\usepackage{lineno}
\modulolinenumbers[1]
\usepackage{setspace}
\usepackage{threeparttable}
\usepackage{threeparttablex}
\usepackage{tikz}
\usepackage{tabularx}
\usepackage{caption}
\usepackage{textcomp}
\usepackage{hyperref}
\usepackage{graphicx}
\usepackage[section]{placeins}
\usepackage{float}

\captionsetup[table]{font=small,labelfont=bf}
\newcommand{\elsvtabsetup}{
  \footnotesize
  \setlength{\tabcolsep}{6pt}
  \renewcommand{\arraystretch}{1.15}
}
\sisetup{
  detect-all      = true,
  round-mode      = places,
  round-precision = 2
}

\setcounter{topnumber}{2}
\setcounter{bottomnumber}{0}
\setcounter{totalnumber}{4}
\setcounter{dbltopnumber}{2}
\renewcommand{\topfraction}{0.95}
\renewcommand{\dbltopfraction}{0.95}
\renewcommand{\bottomfraction}{0}
\renewcommand{\textfraction}{0.05}
\renewcommand{\floatpagefraction}{0.80}
\renewcommand{\dblfloatpagefraction}{0.80}

\hypersetup{
  hidelinks,
  pdftitle = {HarmoSSM for Bi-Temporal SAR Flood Change Detection with Change-Defined Depth Estimation},
  pdfauthor = {Yukun Fan, Zhongqi Shi, Hong Zhang, Yuzhou Liu, Mingyang Song, Yazhe Xie, Yixian Tang}
}

\makeatletter
\setlength{\@fptop}{0pt}
\let\ps@pprintTitle\ps@plain
\makeatother
\setlength{\textfloatsep}{8pt plus 2pt minus 2pt}

\newcommand{\paperlang}[2]{#1}
\journal{Remote Sensing of Environment}

\begin{document}
\doublespacing
\pagenumbering{arabic}
\pagestyle{plain}

\begin{frontmatter}

\title{HarmoSSM for Bi-Temporal SAR Flood Change Detection with Change-Defined Depth Estimation}

\author[inst1,inst2,inst3]{Yukun Fan}
\author[inst4,inst5]{Zhongqi Shi}
\author[inst2,inst1,inst3]{Hong Zhang\corref{cor1}}
\ead{zhanghong@radi.ac.cn}
\author[inst4,inst5]{Yuzhou Liu}
\author[inst1,inst2,inst3]{Mingyang Song}
\author[inst1,inst2,inst3]{Yazhe Xie}
\author[inst1,inst2,inst3]{Yixian Tang}

\cortext[cor1]{Corresponding author}

\affiliation[inst1]{
  organization={Key Laboratory of Digital Earth Science, Aerospace Information Research Institute, Chinese Academy of Sciences},
  city={Beijing},
  postcode={100094},
  country={China}
}

\affiliation[inst2]{
  organization={International Research Center of Big Data for Sustainable Development Goals},
  city={Beijing},
  postcode={100094},
  country={China}
}

\affiliation[inst3]{
  organization={University of Chinese Academy of Sciences},
  city={Beijing},
  postcode={100049},
  country={China}
}

\affiliation[inst4]{
  organization={Shenzhen Technology Institute of Urban Public Safety},
  city={Shenzhen},
  postcode={518000},
  country={China}
}

\affiliation[inst5]{
  organization={Urban Safety Development Institute of Science and Technology (Shenzhen)},
  city={Shenzhen},
  state={Guangdong},
  postcode={518023},
  country={China}
}

\begin{abstract}
Rapid flood assessment requires information on both inundation extent and water depth. Synthetic aperture radar (SAR) supports observation during adverse weather, but complex surface conditions complicate temporal feature comparison, while extensive inundation may coexist with small or weak-contrast patches, narrow corridors, and fragmented margins. Moreover, change-defined inundation excludes pre-existing water, so not all its boundaries provide dry-shoreline elevation constraints. We develop a two-stage flood mapping approach comprising HarmoSSM change detection and Change-defined Flood Depth Estimation (CFDepth). HarmoSSM combines paired states and differences through harmonized change encoding, then uses omni-scale state-space decoding to reconstruct newly flooded extents from regional context and shallow detail. CFDepth combines this extent, screened boundary elevations, and a digital elevation model (DEM) to estimate first-order depth within the supplied flood support. Reference panels from Zhengzhou and Zhuozhou, covering built-up land, cropland, and river-adjacent terrain, illustrate localized and narrow inundation alongside broad flooded areas and fragmented margins. In a separate comparison of existing Brazos depth products, positive-depth coverage over the positive-depth reference domain is 95.81\% for the product labeled CFDepth and 41.91\% for the Floodwater Depth Estimation Tool (FwDET). Mean absolute errors on common-valid pixels are 1.016 and 1.673 m, respectively, although both products have low spatial correlation and negative Nash--Sutcliffe efficiency. The approach combines flood-change mapping with depth estimation conditional on the supplied extent and terrain, providing complementary information on inundation distribution and severity.
\end{abstract}

\begin{keyword}
SAR flood change detection \sep Bi-temporal remote sensing \sep Complex landscapes \sep Multiscale representation \sep DEM-constrained flood depth
\end{keyword}

\end{frontmatter}
\linenumbers

\section{Introduction}

Floods are among the most frequent and widespread natural disasters worldwide, causing thousands of deaths and tens of billions of US dollars in economic losses each year and posing a severe threat to human life and socioeconomic activities \citep{credUndrr2020disasters}. Satellite observations further indicate an upward trend in the proportion of the population exposed within observed flood footprints \citep{tellman2021exposure}, underscoring the urgency of obtaining timely information on the spatial distribution of flood impacts. Timely flood observations over extensive areas can support emergency response, impact assessment, and post-disaster recovery. Through its active microwave imaging capability, synthetic aperture radar (SAR) can acquire surface information regardless of cloud cover and daylight conditions and has become a key data source for flood monitoring \citep{wagner2026gfm}. Reliably extracting spatial information on flood impacts from SAR observations therefore provides an important basis for flood response and assessment.

Flood extent and water depth provide complementary flood information: the former identifies affected areas, while the latter describes the degree of submergence within them. Together, they provide key parameters for assessing risks to people and assets; for example, flood risk assessment for people and vehicles in Zhengzhou combines inundation conditions with hydraulic instability to determine hazard levels \citep{dong2022urbanrisk}. However, obtaining these two types of information from SAR involves different observational constraints. Specular reflection from smooth open water generally weakens the backscattered signal, providing a basis for distinguishing inundation from land \citep{amitrano2024sarFloodReview}. Moreover, backscatter itself does not directly indicate water depth; combining the mapped inundation extent with terrain elevations offers an effective approach to estimating water-surface elevation and inundation depth \citep{cian2018flood}. To address these information needs, combining SAR-based inundation delineation with terrain-assisted depth estimation can provide complementary information on the horizontal distribution of affected areas and the degree of submergence within them.

Flood change detection must distinguish pre-existing water from flood-induced changes while accounting for the diversity of inundation responses across different surface backgrounds. Differences in pre-event backscatter levels and land-cover conditions lead to variations in flood-induced backscatter changes: among inundated areas exhibiting reduced backscatter, some show pronounced changes, whereas others produce only weak responses, complicating the extraction and discrimination of flood-change features \citep{giustarini2013urbanFloodCd,amitrano2024sarFloodReview}. To support flood discrimination in complex settings, DAM-Net combines change features with water-change semantics to enrich flood-change representations \citep{saleh2024dam}, while AWCA-Net integrates neighborhood context and difference-guided attention to enhance the discrimination of subtle and heterogeneous changes \citep{du2026high}. In addition, a global SAR flood-mapping study uses auxiliary land-cover data to remove permanent water, providing supplementary constraints for identifying event-related inundation \citep{misra2025globalFloods}. These studies enrich the evidence available for flood discrimination across complex surface backgrounds through feature fusion, context modeling, and auxiliary information constraints. However, even within inundated areas, the magnitude of backscatter changes can vary with the surface background; decisions that rely too heavily on pronounced changes may therefore omit inundated areas with weaker responses. A change-discrimination mechanism that accommodates these variations in inundation response across different surface backgrounds is therefore needed to improve the recognition of weak inundation changes alongside pronounced ones, enabling more complete identification of inundated areas.

The spatial delineation of flooding must accommodate extensive contiguous inundation, narrow flooded strips, and scattered patches, requiring extent reconstruction to use both information about the affected region as a whole and local positioning cues near its boundaries. SAR flood-mapping studies have explored multiscale representation, selective feature fusion, and surface-context constraints to address this need. MS-Deeplab combines features at different levels for water extraction \citep{wu2022msdeeplab}, Siam-DWENet introduces a multiscale pyramid into inundation detection from paired images \citep{zhao2023siamdwenet}, and SISCNet extracts features from different neighborhoods for bi-temporal inundation discrimination \citep{tahermanesh2025siscnet}. For feature selection, attentive fusion weights spatial and channel information \citep{yadav2022attentive}, while adaptive-window methods organize multiscale features using local neighborhoods \citep{du2026high}; land cover, historical water occurrence, and hydrological information also supplement image-based interpretation \citep{zhu2025globalfloodsar}. In related change-detection research, BIT organizes regional information through bi-temporal semantic interaction \citep{chen2021remote}, while CGNet uses deep change priors to guide multiscale feature fusion \citep{han2023change}. These studies enrich the multiscale information, regional context, and local spatial cues available for flood-extent identification. However, narrow flooded strips and scattered inundation may be omitted if local positioning cues are insufficiently preserved during regional information aggregation. Conversely, decisions that rely too heavily on local responses may lack the regional context needed to maintain consistent predictions within contiguous inundation. A reconstruction mechanism that connects regional context aggregation with local-detail recovery is therefore needed to preserve the completeness of contiguous flood coverage while accurately delineating narrow, scattered, and fragmented flooded areas.

Identifying the spatial distribution of affected areas alone is insufficient to support disaster assessment, and water-depth information is also needed to characterize the degree of submergence. Methods based on flood observations and terrain commonly use ground elevations at wet--dry boundaries to constrain water-surface elevation and calculate depth as the difference between water-surface and ground elevations. Cian et al. estimated water levels from statistics of lidar terrain elevations along the edges of SAR-identified areas \citep{cian2018flood}, while FwDET v2.0 improved boundary identification and local water-level assignment to reduce inappropriate allocation of boundary elevations across permanent water bodies \citep{cohen2019floodwater}. For water-surface inference, FLEXTH combines boundary elevations with terrain conditions to estimate water levels \citep{betterle2024water}, and RS-FloodXDepth integrates boundary-elevation analysis with hydrologically guided spatial inference \citep{tian2026floodxdepth}. Yang et al. connected water segmentation with FLEXTH-based geometric depth estimation, illustrating a way to link image-based identification with terrain-based depth estimation \citep{yang2026csmamba}. These studies provide a methodological basis for converting spatial flood information into water-depth information and indicate that depth estimation depends both on the selection of boundary elevations and on how these local references are used to infer the interior water surface.

However, when the boundary-elevation-based depth-estimation methods discussed above are applied directly to change-detection results, not all detected boundaries meet the conditions for inferring water levels from ground elevations. Change detection identifies newly flooded areas whose boundaries may adjoin permanent water bodies, leaving areas classified as unchanged outside these boundaries still water-covered, so such boundaries cannot directly serve as wet--dry boundaries that provide water-level references \citep{betterle2024water}. Even along boundaries that satisfy wet--dry shoreline conditions, locally steep slopes and anomalous elevations may introduce biases into water-level references. Directly assigning these elevations to interior pixels may propagate the biases into the interior of the region \citep{cohen2019floodwater}. When adjacent pixels use boundary references with different elevations, this assignment may also introduce local jumps in the water surface. Depth estimation therefore requires screening boundary elevations to obtain stable constraints, provided that the boundaries meet the conditions for their use, and coordinating water-surface elevations at adjacent locations while preserving these constraints to reduce local bias propagation and discontinuities caused by direct assignment.

In summary, flood change detection focuses on identifying flooded areas, whereas depth estimation relies on inferring the water surface under terrain constraints, and the two tasks remain insufficiently connected in translating detection results into effective constraints for depth estimation. Yang et al. integrated water identification and depth estimation, illustrating the practical value of combining the two tasks for flood mapping \citep{yang2026csmamba}. Combining change detection with depth estimation therefore requires both accurate identification of newly flooded areas and the establishment of appropriate elevation constraints, enabling comprehensive flood mapping that captures both horizontal distribution and the vertical degree of submergence and provides more complete spatial information for disaster assessment.

To address the above problems, this study proposes a two-stage mapping approach that combines flood change detection with depth estimation. In the change-detection stage, HarmoSSM harmonizes bi-temporal representations to address variations in change-response strength and combines regional context with local detail to delineate both contiguous inundation and narrow, scattered flooded areas. In the depth-estimation stage, CFDepth combines detection results with terrain information, using boundary-elevation screening and interior water-surface coordination to address local elevation biases and discontinuities caused by direct assignment. The approach links the identification of newly flooded areas with depth estimation within them in a unified mapping workflow, providing a comprehensive characterization of flood distribution and the degree of submergence.

The main contributions of this study are as follows:
\begin{enumerate}
  \item \textbf{Harmonized change encoding that jointly uses state and difference information.} A Harmonized Change Encoder (HCE) is designed to compare pre-event and post-event features against a common statistical reference, jointly encode the states at both times and their differences, and aggregate regional information through query interaction, providing a basis for discriminating changes of varying response strengths using both state information and change cues.
  \item \textbf{Omni-scale change decoding combining regional aggregation and detail recovery.} An Omni-Scale State-Space Change Decoder (OSCD) is designed to aggregate regional context from deep change features and progressively incorporate shallow spatial detail, reconstructing multiscale change information into flood extents that capture both contiguous coverage and the locations of narrow, scattered flooded areas.
  \item \textbf{Depth estimation combining boundary constraints with water-surface coordination.} CFDepth is developed to screen boundary elevations based on slope and local elevation consistency, establish and preserve water-surface elevation constraints, and refine the continuity of the interior water surface to estimate first-order water depth within the supplied detection area under the available terrain conditions.
\end{enumerate}
\clearpage

\section{Study Areas and Data}

\subsection{Study Areas and Flood Settings}

Three study areas are selected: central Zhengzhou in Henan Province, Zhuozhou in Hebei Province, and an area near Binyang County and Hengzhou in Guangxi. Their locations and land cover are shown in Fig.~\ref{fig:dataset}. The three settings comprise concentrated built-up land, interspersed cropland and settlements, and cropland and settlements distributed among tree-covered areas, respectively. They provide observations for examining flood-change detection under different surface conditions and cases for terrain-constrained depth estimation within the detected regions.

The Zhengzhou study area covers the city centre and corresponds to the July 2021 flood. Built-up land is concentrated in Fig.~\ref{fig:dataset}(b), with channels and other water bodies crossing or adjoining developed areas and cropland and vegetation along the periphery. Flood photographs show standing water around roads, buildings, and vehicles, illustrating the complexity of locally inundated spaces among urban features. This setting is used to examine the identification of scattered inundation and complex boundaries in a built environment.

The Zhuozhou study area corresponds to the flood of late July to early August 2023, for which the Haihe basin study by \citet{lan2025haihe} provides regional context. Cropland and settlements are interspersed with water bodies in Fig.~\ref{fig:dataset}(c), and the photographs show inundation around roads, transport infrastructure, and buildings. Compared with central Zhengzhou, this area provides a setting where agricultural and built environments adjoin, allowing examination of the spatial continuity and local boundaries of inundated regions across different surfaces.

The Guangxi study area lies near Binyang County and Hengzhou and is delineated by the common valid coverage of the selected bi-temporal LT-1 imagery. The administrative names indicate its location rather than its boundary. Persistent heavy rainfall affected Nanning during 4--6 July 2026 under the influence of Typhoon Maysak. On the morning of 6 July, overtopping and breaches were reported at the Liulan and Yunbiao reservoirs in Hengzhou, while the Liuwang reservoir in Binyang County overtopped, with flooding in some downstream areas. Figure~\ref{fig:dataset}(d) shows cropland, settlements, and river waters distributed among tree-covered areas, while the photographs illustrate inundation around settlements and roads. This setting complements the other two areas in both land-cover pattern and SAR observation conditions for examining flood-change detection and its connection to terrain-constrained depth estimation.

\begin{figure}[H]
  \centering
  \includegraphics[width=\textwidth,height=.65\textheight,keepaspectratio]{figure/figure1.jpg}
  \caption{Locations, land cover, and flood settings of the three study areas. Panel (a) shows their locations in China, and (b)--(d) show land cover in central Zhengzhou, Zhuozhou, and Guangxi, respectively. P1--P4, P5--P8, and P9--P12 illustrate local flood conditions in the three areas.}
  \label{fig:dataset}
\end{figure}

\subsection{Remote Sensing Observations and Terrain Data}\label{sec:depth_data}

Bi-temporal SAR imagery provides reference states and flood-period observations for the three areas, with key parameters summarized in Table~\ref{tab:gf3_eval}. Zhengzhou and Zhuozhou use HV-polarized GF3 imagery. Their reference and flood-period observations were acquired in July but one year apart, so surface changes over the intervening period still need to be distinguished from flood-related changes. Guangxi uses HH-polarized LT-1B imagery, with two scenes mosaicked for each acquisition date.

\begin{table*}[htbp]
\centering
\captionsetup{name=Table}
\caption{SAR observations and terrain data.}
\label{tab:gf3_eval}
\small
\begin{tabular}{lccccc}
\toprule
\multicolumn{6}{l}{(a) GF3 observations} \\
\midrule
Scene & Mode & Pixel spacing & Polarization & Pre-event & Post-event \\
\midrule
Zhengzhou & FSII & 10 m & HV & 2020-07-20 & 2021-07-20 \\
Zhuozhou & FSII & 10 m & HV & 2022-07-20 & 2023-07-31 \\
\midrule
\multicolumn{6}{l}{(b) Terrain data} \\
\midrule
Product & \multicolumn{2}{c}{Sampling} & \multicolumn{3}{l}{Role and source} \\
FABDEM & \multicolumn{2}{c}{Approximately 30 m} & \multicolumn{3}{l}{Terrain input for depth estimation} \\
 & \multicolumn{2}{c}{} & \multicolumn{3}{l}{\citet{hawker2022fabdem}} \\
\bottomrule
\end{tabular}
\end{table*}


All three study areas use FABDEM as the terrain data source, with a sampling interval of approximately 30 m. The product reduces the effects of building and forest heights on elevation to better represent ground elevations \citep{hawker2022fabdem}. FABDEM provides the ground-elevation reference and water-surface elevation constraints at valid boundaries for CFDepth, supporting first-order depth estimation within the supplied detection regions. The spatial detail of the estimates is limited by terrain sampling and elevation errors.

\section{Methodology}\label{sec:methodology}

\subsection{Overview}\label{sec:overview}

Flood mapping under complex surface conditions requires information on both the location of newly affected areas and the degree of inundation within them. To meet these needs, we develop a two-stage framework comprising bi-temporal SAR change detection and depth estimation within the detected extent (Fig.~\ref{fig:framework}). Registered pre-event and post-event SAR images, denoted by \(\mathbf{I}^A\) and \(\mathbf{I}^B\), provide the observations for identifying newly flooded areas, while the DEM supplies terrain elevations for depth estimation. The first stage uses HarmoSSM to produce the newly flooded extent \(\hat{\mathbf{M}}\). The second uses CFDepth to estimate internal depth from this extent and the DEM, spatially linking extent identification with the estimation of inundation severity.

The change detection stage uses HarmoSSM, comprising a Harmonized Change Encoder (HCE) and an Omni-Scale State-Space Change Decoder (OSCD). HCE organizes shared spatial representations, semantic information, paired harmonization, and query interaction into a continuous change encoding process, yielding five levels of change features. OSCD aggregates deep regional context and then progressively incorporates shallow local detail to reconstruct pixel-level predictions of new inundation. The detection target is binary, with \(y=1\) denoting new inundation and \(y=0\) denoting background. Softmax converts the two-class logits into foreground probabilities, and the selected threshold produces \(\hat{\mathbf{M}}\) over valid input pixels as the spatial input to the depth estimation stage.

In the depth estimation stage, CFDepth first intersects the detected extent with valid DEM coverage, screens boundary locations using slope and local elevation, and constructs water-surface anchors from terrain elevations at the boundary and its inner neighborhood. Anchor elevations propagate along reachable paths within the extent to initialize the water surface. Local continuity refinement then retains the anchors and enforces the terrain lower bound. Subtracting the DEM from the resulting water-surface elevation (WSE) yields the depth distribution within the given extent. The two stages provide horizontal flood distribution and vertical information, with depth estimates conditioned on the detected extent, valid terrain, and selected boundary constraints.

\begin{figure}[H]
  \centering
  \includegraphics[width=\linewidth,height=0.70\textheight,keepaspectratio]{figure/figure2.jpg}
  \caption{Overall workflow of HarmoSSM and CFDepth.}
  \label{fig:framework}
\end{figure}

\subsection{Harmonized Change Encoder (HCE)}\label{sec:encoder}\label{sec:ha}\label{sec:cqi}

New inundation discrimination requires information about surface states at both times and their changes, together with a comparable representation that relates them. The Harmonized Change Encoder (HCE) organizes shared representation extraction, paired harmonization, and change query interaction into a continuous process. It first extracts multiscale features for each time, then addresses the statistical reference and local correspondence, and finally combines state and difference information into change representations. HCE thus outputs five levels of change features for the decoder, expressing bi-temporal comparison at different spatial scales and supplying the inputs for flood-extent reconstruction.

Shared spatial representations retain location cues, while semantic assistance supplements the interpretation of surfaces within their surroundings. A shared convolutional encoder and feature pyramid network extract spatial representations \(P1\)--\(P5\) \citep{tan2019efficientnet,lin2017fpn}, with progressively lower resolutions. Shallow levels retain finer position and texture cues, whereas deeper levels organize information over larger areas. Frozen DINOv3 representations are adapted and incorporated into \(P3\)--\(P5\) \citep{simeoni2025dinov3}, where they are concatenated with the corresponding CNN features and fused through convolution and CBAM \citep{woo2018cbam}; \(P1/P2\) retain convolutional features. The upper part of Fig.~\ref{fig:encoder_ha} shows these scale-dependent paths. The CNN and semantic branches each share their respective weights across the pre-event and post-event inputs, while fusion is performed separately within each time. Both times therefore use the same feature mappings while retaining their individual state information. The output at level \(i\) and time \(t\) is denoted by \(\mathbf P_i^t\).

\begin{figure}[H]
  \centering
  \includegraphics[width=\linewidth,height=0.70\textheight,keepaspectratio]{figure/figure3.jpg}
  \caption{Shared representation and paired harmonization in HCE. The upper part shows the bi-temporal shared CNN--FPN, DINO semantic fusion, and \(P1\)--\(P5\) paths; the lower part details paired statistical calibration and local soft alignment. \(A/B\) denote pre-event/post-event inputs; \(Q\) comes from post-event features and sampled \(K/V\) from pre-event features. Local alignment updates the pre-event branch.}
  \label{fig:encoder_ha}
\end{figure}

Paired statistical calibration establishes a common reference for shallow-feature comparison from the current image pair. At \(P1/P2\), Harmonized Alignment (HA) computes channel-wise means \(\boldsymbol\mu_i^{AB}\) and variances \(\mathbf v_i^{AB}\) jointly over both times and spatial locations. The same statistics normalize both feature sets before a shared learnable affine transformation:
\begin{equation}\label{eq:ha_calibration}
 \mathbf{G}_i^t=\frac{\mathbf{P}_i^t-\boldsymbol\mu_i^{AB}}{\sqrt{\mathbf{v}_i^{AB}+\epsilon}}
 \odot\left[\operatorname{softplus}(\mathbf{s}_i^c)+\epsilon\right]+\boldsymbol\mu_i^c,
 \qquad \mathbf{U}_i^t=\mathbf{G}_i^t+\mathcal{R}_i(\mathbf{G}_i^t).
\end{equation}
Here, \(t\in\{A,B\}\) denotes the pre-event and post-event times, respectively, \(i\in\{1,2\}\), and \(\epsilon\) is a numerical stabilizer. The affine location and scale parameters \(\boldsymbol\mu_i^c\) and \(\mathbf s_i^c\) are learned during training. Pair-dependent statistics give both feature sets the same normalization reference, while the affine parameters learn the scale and offset of this common space. The convolutional mapping \(\mathcal R_i\) takes the transformed feature \(\mathbf G_i^t\) as input and adds its residual output to \(\mathbf G_i^t\). This provides a learnable recovery path for the calibrated local representation, yielding \(\mathbf U_i^t\).

Local correspondence processing then relates spatial features across the two times within a limited neighborhood. Statistical calibration acts on their common reference; local soft alignment uses post-event features as queries and samples keys and values from pre-event features to construct a local correction to the pre-event representation. At \(P1/P2\), calibration precedes alignment, with \(\mathbf B_i^t=\mathbf U_i^t\); at \(P3\), alignment directly receives the semantically fused features \(\mathbf B_3^t=\mathbf P_3^t\). The pre-event update is
\begin{equation}\label{eq:ha_alignment}
  \tilde{\mathbf{P}}_{i}^{A}(\mathbf{p})
  =
  \mathbf{B}_{i}^{A}(\mathbf{p})
  +
  \mathcal{O}_{i}\!\left[
  \sum_{k=1}^{K}
  a_{ik}(\mathbf{p})\,
  \mathbf{V}_{i}^{A}
  \!\left(\mathbf{p}+\mathbf{r}_{k}+\Delta\mathbf{o}_{ik}(\mathbf{p})\right)
  \right],
  \quad i\in\{1,2,3\}.
\end{equation}
Here, \(\mathbf p\) is a grid position, \(\mathbf r_k\) is the \(k\)th reference sampling offset, \(\Delta\mathbf o_{ik}\) is a learned local offset, and \(\mathbf V_i^A\) is the pre-event value mapping. Paired features and their directed difference jointly predict offsets and attention biases. Offsets change where pre-event keys and values are sampled; biases enter the matching scores between post-event queries and sampled keys before local softmax produces \(a_{ik}\). The mapping \(\mathcal O_i\) combines weighted values across heads and updates the pre-event feature through a residual. The local-alignment stage leaves its post-event input \(\mathbf B_i^B\) unchanged, and \(P4/P5\) retain their incoming features. The resulting pyramids \(\{\mathbf X_i^A,\mathbf X_i^B\}_{i=1}^5\) supply subsequent interaction. These paths connect shallow statistical harmonization with local spatial correspondence while defining the processing applied at each scale.

Change Query Interaction (CQI) jointly organizes state and difference information from the harmonized paired representations (Fig.~\ref{fig:cqi}). Pre-event and post-event features retain the surface states before and after a change; the directed difference preserves the direction of feature change, and the absolute difference supplies a sign-independent magnitude cue. At each scale, convolutional projection \(\mathcal P_i\) combines these four inputs into the dense pair feature \(\mathbf Z_i\):
\begin{equation}\label{eq:cqi_pair}
  \mathbf{Z}_{i}
  =
  \mathcal{P}_{i}\!\left(
  \left[
  \mathbf{X}_{i}^{A},
  \mathbf{X}_{i}^{B},
  \mathbf{X}_{i}^{B}-\mathbf{X}_{i}^{A},
  \left|\mathbf{X}_{i}^{B}-\mathbf{X}_{i}^{A}\right|
  \right]\right),
  \quad i=1,\ldots,5,
\end{equation}
Brackets denote channel concatenation. State features provide the before-and-after context for interpreting differences, while difference features explicitly express the relationship between the two times. These quantities belong to the learned representation and do not directly represent backscatter differences or water depth. Flattening \(\mathbf Z_i\) into spatial tokens \(\mathbf T_i=\operatorname{Flatten}(\mathbf Z_i)\) preserves the correspondence between each token and its grid position, providing a common dense basis for information aggregation and spatial restoration.

Query interaction uses two steps, from spatial information to queries and from queries back to spatial positions, to incorporate regional relationships into dense change encoding. At each scale, learnable queries \(\mathbf Q_i^0\) first aggregate information from \(\mathbf T_i\). The updated queries \(\mathbf Q_i^1\) then serve as keys and values to return context to the original tokens. The two multi-head attention (MHA) operations therefore have different information directions:
\begin{align}
 \mathbf{Q}_{i}^{1}
 &=\operatorname{LN}_{q,i}\!\left[\mathbf{Q}_{i}^{0}
 +\operatorname{MHA}_{q\leftarrow t}(\mathbf{Q}_{i}^{0},\mathbf{T}_{i},\mathbf{T}_{i})\right],\label{eq:cqi_query}\\
 \mathbf{T}_{i}^{1}
 &=\operatorname{LN}_{t,i}\!\left[\mathbf{T}_{i}
 +\operatorname{MHA}_{t\leftarrow q}(\mathbf{T}_{i},\mathbf{Q}_{i}^{1},\mathbf{Q}_{i}^{1})\right].\label{eq:cqi_token}
\end{align}
The operators \(\operatorname{LN}_{q,i}\) and \(\operatorname{LN}_{t,i}\) denote separate query and token layer normalization. The first step gathers spatial information into a limited set of queries, while the second retains the original tokens as query inputs so that returned information corresponds to their positions. Both attention updates retain input residuals. A residual feed-forward network (FFN) then updates the tokens and reuses \(\operatorname{LN}_{t,i}\). After spatial reshaping, the features are added to the original dense feature \(\mathbf Z_i\) and passed through output convolution and CBAM:
\begin{equation}\label{eq:cqi_output}
 \mathbf{D}_i=\mathcal{O}_i\!\left[
 \operatorname{Reshape}\!\left(\operatorname{LN}_{t,i}\left[\mathbf{T}_{i}^{1}+\operatorname{FFN}_i(\mathbf{T}_{i}^{1})\right]\right)+\mathbf{Z}_i\right].
\end{equation}
Here, \(\mathcal O_i\) denotes the CQI output mapping. The dense residual preserves a direct path for the original paired information alongside query interaction. Each of the five scales has independent query and interaction parameters, exchanges information within that scale, and outputs \(\mathbf D_1\)--\(\mathbf D_5\) for OSCD flood-extent reconstruction. Input conversion and numerical settings are specified in Section~\ref{sec:eval_protocol}.

\begin{figure}[H]
  \centering
  \includegraphics[width=\linewidth,height=0.70\textheight,keepaspectratio]{figure/figure4.jpg}
  \caption{Change query interaction in HCE. Paired states, directed difference, and absolute difference are projected into dense feature \(Z_i\), flattened into \(T_i\), and processed by query and spatial-token updates. A residual FFN and dense residual then produce change output \(D_i\). \(Q_i^0/Q_i^1\) denote queries before/after updating; \(C\) and plus signs denote concatenation and addition.}
  \label{fig:cqi}
\end{figure}

\subsection{Omni-Scale State-Space Change Decoder (OSCD)}\label{sec:oscd}\label{sec:loss}

Flood-extent reconstruction must translate regional change judgments into the spatial coverage and local positions of affected areas. The Omni-Scale State-Space Change Decoder (OSCD) organizes deep context aggregation followed by shallow-detail recovery. Features \(\mathbf D_3\)--\(\mathbf D_5\) establish a regional representation \(\mathbf R_3\), which is then combined with \(\mathbf D_2\) and \(\mathbf D_1\) on progressively finer grids. The left part of Fig.~\ref{fig:oscd} shows the complete path from HCE change features to extent prediction, connecting broader change relationships with local positioning.

Deep fusion first establishes a common spatial grid for change information from different scales. Features \(\mathbf D_3\)--\(\mathbf D_5\) undergo separate channel projections \(\rho_i^c\) and are resized to the \(\mathbf D_3\) grid, producing \(\mathbf E_3,\mathbf E_4,\mathbf E_5\). Concatenation and fusion then yield \(\mathbf F\):
\begin{equation}\label{eq:oscd_deep_fusion}
 \mathbf E_i=\operatorname{Resize}_{3}\!\left(\rho_i^c(\mathbf D_i)\right),
 \quad i\in\{3,4,5\},\qquad
 \mathbf F=\operatorname{Fuse}\!\left[\mathbf E_3,\mathbf E_4,\mathbf E_5\right].
\end{equation}
Here, \(\operatorname{Resize}_3\) adjusts features to the third-level spatial dimensions, and \(\operatorname{Fuse}\) denotes convolutional fusion after concatenation. The common grid permits position-wise combination of information from the three levels, while the separately retained \(\mathbf E_i\) also support subsequent context reinjection. This context structure follows the multiscale organization of SegMAN \citep{fu2025segman}; here it receives deep change features from HCE for bi-temporal flood-extent reconstruction.

Multiresolution input construction further organizes information with different spatial coverage from the fused feature, as shown in the lower-right part of Fig.~\ref{fig:oscd}. Feature \(\mathbf F\) follows original-resolution, twofold-downsampling, and fourfold-downsampling paths. Pixel unshuffle rearranges spatial positions into channels in the first two paths, bringing all three to the same coarser grid. Separate channel reductions and concatenation produce the scan input \(\mathbf G\):
\begin{equation}\label{eq:oscd_scan_input}
 \mathbf G=\left[
 \rho_0\!\left(\operatorname{PU}_{4}(\mathbf F)\right),
 \rho_1\!\left(\operatorname{PU}_{2}(\mathcal C_{2}(\mathbf F))\right),
 \rho_2\!\left(\mathcal C_{4}(\mathbf F)\right)
 \right].
\end{equation}
Here, \(\mathcal C_2\) and \(\mathcal C_4\) denote stride-2 and stride-4 convolutions, \(\operatorname{PU}_r\) rearranges each \(r\times r\) spatial block into channels, and \(\rho_0,\rho_1,\rho_2\) reduce channels in the respective paths. Spatial rearrangement expresses the local arrangement at the original resolution and the downsampled regional information on a common grid for a shared context module. The equation omits boundary padding where needed; dimensions, convolution kernels, and padding rules are specified in Section~\ref{sec:eval_protocol}.

State-space context aggregates regional information on the common grid and returns it to the deep change representations. The operator \(\operatorname{SSContext}\) unfolds sequences in horizontal, vertical, and their reverse directions, applies selective state-space scanning, and restores and merges the results on the spatial grid \citep{gu2023mamba,liu2024vmamba}, yielding \(\mathbf O\). Upsampling \(\mathbf O\) to the context grid produces \(\mathbf S\), which is then projected by \(\rho^r\) into \(\mathbf J\) with channels matching \(\mathbf E_i\). The regional representation is constructed as
\begin{equation}\label{eq:oscd_context_flow}
 \mathbf O=\operatorname{SSContext}(\mathbf G),\qquad
 \mathbf S=\operatorname{Up}_{3}(\mathbf O),\qquad
 \mathbf J=\rho^r(\mathbf S).
\end{equation}
\begin{equation}\label{eq:oscd_context_reconstruction}
 \begin{aligned}
 \mathbf R_3=\operatorname{Reconstruct}\!\big[&\mathbf S,\operatorname{Short}(\mathbf F),
 \operatorname{Global}_{3}(\mathbf F),\\
 &\mathbf E_3+\mathbf J,\mathbf E_4+\mathbf J,\mathbf E_5+\mathbf J\big].
 \end{aligned}
\end{equation}
The operator \(\operatorname{Short}\) is a short convolutional branch applied directly to \(\mathbf F\), while \(\operatorname{Global}_3\) includes global pooling, projection, and resizing to the context grid. Upsampled features \(\mathbf S\) enter concatenation directly, and \(\mathbf J\) is added separately to each \(\mathbf E_i\), supplying a common context contribution to the deep levels. The upper-right part of Fig.~\ref{fig:oscd} corresponds to this six-path aggregation: the scanned representation, local convolution, global information, and three reinjected deep features are concatenated and reconstructed into \(\mathbf R_3\). Context return and retention of the original level-specific information are thus organized within the same aggregation structure.

Shallow-detail recovery carries \(\mathbf R_3\) toward higher resolutions while introducing the corresponding local position cues at each level. First, \(\mathbf R_3\) is upsampled to the \(\mathbf D_2\) grid and concatenated with projected \(\mathbf D_2\) to reconstruct \(\mathbf R_2\). The same procedure incorporates \(\mathbf D_1\) to obtain \(\mathbf R_1\):
\begin{equation}\label{eq:oscd_reconstruction}
  \mathbf{R}_{2}
  =
  \mathcal{R}_{2}\!\left[
  \operatorname{Up}(\mathbf{R}_{3}),\rho_{2}^{d}(\mathbf{D}_{2})
  \right],
  \qquad
  \mathbf{R}_{1}
  =
  \mathcal{R}_{1}\!\left[
  \operatorname{Up}(\mathbf{R}_{2}),\rho_{1}^{d}(\mathbf{D}_{1})
  \right].
\end{equation}
Here, \(\rho_i^d\) is a detail projection, \(\mathcal R_i\) is the convolutional reconstruction mapping at that level, and \(\operatorname{Up}\) resizes to the current shallow feature dimensions. Concatenation presents the aggregated regional information and level-specific detail jointly to the reconstruction mapping, progressively locating the extent judgment on finer grids. A convolutional head maps \(\mathbf R_1\) to two-class logits, which are bilinearly resized to the input grid; softmax and the selected threshold then convert them into the newly flooded extent.

\begin{figure}[H]
  \centering
  \includegraphics[width=\linewidth,height=0.70\textheight,keepaspectratio]{figure/figure5.jpg}
  \caption{Flood-extent reconstruction in OSCD. The left part shows regional aggregation of \(D_3\)--\(D_5\), progressive recovery with \(D_2/D_1\), and extent output. The upper-right part details regional-context aggregation, and the lower-right part constructs multiresolution input \(G\) from fused feature \(F\). \(O\), \(S\), and \(J\) denote context output, its upsampled representation, and the reinjection projection; \(R_i\) denotes reconstructed features and PU denotes pixel unshuffle.}
  \label{fig:oscd}
\end{figure}

The training objective constrains both the final extent prediction and the five HCE change-feature levels, jointly supervising reconstruction and multiscale change encoding. Five auxiliary heads directly receive \(\mathbf D_1\)--\(\mathbf D_5\), and the main and auxiliary predictions contribute to pixel-classification and foreground-overlap losses. Shallow auxiliary predictions provide a detached support reference. A support-preservation term constrains the main response at supported locations, while coarse consistency constrains excessive responses where support is low. The total objective is
\begin{equation}\label{eq:total_loss}
  \mathcal{L}
  =
  \mathcal{L}_{\mathrm{seg}}
  +
  \lambda_{\mathrm{sup}}(e)\mathcal{L}_{\mathrm{sup}}
  +
  \lambda_{\mathrm{coarse}}(e)\mathcal{L}_{\mathrm{coarse}}.
\end{equation}
Here, \(\mathcal L_{\mathrm{seg}}\) combines Focal and Tversky losses for the main and auxiliary outputs, \(\mathcal L_{\mathrm{sup}}\) and \(\mathcal L_{\mathrm{coarse}}\) denote support preservation and coarse consistency, and \(e\) is the training epoch. Auxiliary supervision applies task constraints to HCE change features before decoding, while the support-related terms connect shallow predictions to the final or coarse outputs. Definitions, weights, and schedules are specified in Section~\ref{sec:eval_protocol}.

\subsection{First-Order Depth Estimation Constrained by Inundation Extent and Terrain (CFDepth)}\label{sec:cfdepth}

Depth estimation within an inundated extent requires relating the horizontal flood distribution to terrain elevation. CFDepth takes the newly flooded extent \(\hat{\mathbf M}\) and DEM \(Z\) as inputs. Its estimation support is the intersection with the valid terrain domain \(\Omega_Z\), namely \(\Omega_f=\{p:\hat{\mathbf M}(p)=1\}\cap\Omega_Z\). The extent determines where computation is performed, and the DEM constrains water-surface elevation relative to the ground. Dilation and erosion define the immediately adjacent exterior candidate boundary \(\mathcal B_{\mathrm{raw}}\) and interior boundary \(\mathcal B_{\mathrm{in}}\):
\begin{equation}\label{eq:cfdepth_boundary}
 \mathcal B_{\mathrm{raw}}=\operatorname{Dilate}(\Omega_f)\setminus\Omega_f,
 \qquad
 \mathcal B_{\mathrm{in}}=\Omega_f\setminus\operatorname{Erode}(\Omega_f).
\end{equation}
These boundaries supply exterior terrain samples and interior pairing locations, respectively, defining the spatial domains for subsequent screening and water-surface anchor construction.

Candidate-boundary screening is intended to prevent invalid terrain, steep slopes, and local elevation anomalies from directly entering water-surface constraints. Candidates with valid DEM values and slope \(\theta\) no greater than \(\theta_{\max}\) form \(\mathcal B_c\). Within local neighborhood \(\mathcal N_b(p)\), the median candidate elevation \(m(p)\) defines a pointwise absolute deviation \(e(p)\). A second local median \(a(p)\) is then computed over this deviation field to select the valid exterior boundary \(\mathcal B_v\):
\begin{equation}\label{eq:cfdepth_boundary_screen}
 \begin{aligned}
 \mathcal B_c&=\{p\in\mathcal B_{\mathrm{raw}}\cap\Omega_Z:
       \theta(p)\text{ is finite},\ \theta(p)\leq\theta_{\max}\},\\
 m(p)&=\operatorname{median}_{q\in\mathcal N_b(p)\cap\mathcal B_c} Z(q),
 \qquad e(p)=|Z(p)-m(p)|,\\
 a(p)&=\operatorname{median}_{\substack{q\in\mathcal N_b(p)\cap\mathcal B_c\\e(q)\text{ is finite}}}e(q),\\
 \mathcal B_v&=\{p\in\mathcal B_c:e(p),a(p)\text{ are finite},\
 e(p)\leq\max[\kappa a(p),\tau_{\min}]\}.
 \end{aligned}
\end{equation}
Here, \(\kappa\) is the deviation multiplier and \(\tau_{\min}\) is the minimum elevation tolerance. The second median uses each candidate's deviation from its own local elevation median, preserving the implemented sequence of a local elevation reference, a deviation field, and a local tolerance. The minimum tolerance retains a nonzero acceptance range where deviations are small. This screening constrains terrain-sample validity and local elevation consistency; it does not classify the water-body type adjoining a boundary.

Water-surface anchors pair local terrain samples from the exterior and interior boundaries at interior-boundary locations. For each \(p\in\mathcal B_{\mathrm{in}}\), median elevations within a pairing neighborhood are computed separately from \(\mathcal B_v\) and \(\mathcal B_{\mathrm{in}}\), denoted by \(Z_{\mathrm{dry}}(p)\) and \(Z_{\mathrm{wet}}(p)\). The wet-side statistic uses only interior-boundary samples. Locations where both medians are valid form the anchor set \(\mathcal A\), with water-surface elevation
\begin{equation}\label{eq:cfdepth_anchor}
  S_{b}
  =
  \max\!\left(
  \frac{Z_{\mathrm{wet}}+Z_{\mathrm{dry}}}{2},
  \mathbf{Z}+d_{\min}
  \right).
\end{equation}
All elevations in this equation are evaluated at the same pairing location \(p\in\mathcal A\), and \(d_{\min}\) is the numerical lower bound above terrain. Averaging the two medians supplies a terrain reference at the boundary, and the lower bound keeps the anchor at or above local terrain plus the minimum depth. These anchors are water-surface constraints derived from the detected boundary and terrain; their applicability depends on the supplied extent and boundary samples.

The initial water surface is obtained by propagating the nearest reachable anchor within the support. For each support pixel \(p\), the nearest anchor \(b(p)\) is selected by path distance through eight-connected pixels in \(\Omega_f\), subject to the maximum propagation distance when configured. Where an anchor is available, initialization uses \(S^{(0)}(p)=\max[S_b(b(p)),Z(p)+d_{\min}]\). Where no anchor is reachable or the propagation limit is exceeded, initialization falls back to \(Z(p)+d_{\min}\) and retains a low-confidence flag. Paths within the support determine where anchor information can reach, while the terrain lower bound applies to all initialized pixels. Propagation supplies the starting surface for continuity refinement.

Continuity refinement coordinates neighboring water-surface elevations within the valid support while retaining anchor and terrain constraints. For a support pixel, the local neighborhood mean and current surface are combined using relaxation coefficient \(\eta\):
\begin{equation}\label{eq:cfdepth_continuity}
 S^{(k+1)}(p)=\max\!\left[(1-\eta)S^{(k)}(p)+\eta\operatorname{mean}_{q\in\mathcal{N}(p)\cap\Omega_f}S^{(k)}(q),\ Z(p)+d_{\min}\right].
\end{equation}
The neighborhood \(\mathcal N(p)\) includes the center pixel, and its mean uses only valid water-surface pixels in \(\Omega_f\). Each update applies the terrain lower bound and then restores \(p\in\mathcal A\) to fixed anchor values \(S_b(p)\). Other support pixels participate in refinement, and low-confidence flags from initialization are retained. After \(K\) updates, depth is obtained by subtracting terrain from the final water surface:
\begin{equation}\label{eq:cfdepth_depth}
  \hat{\mathbf{d}}(\mathbf{p})
  =
  \begin{cases}
    \max\!\left(\mathbf{S}^{(K)}(\mathbf{p})-\mathbf{Z}(\mathbf{p}),0\right),
      & \mathbf{p}\in\Omega_f,\\
    \text{NoData}, & \text{otherwise}.
  \end{cases}
\end{equation}
NoData is retained outside the support, keeping the estimation domain consistent with the extent and valid-terrain conditions. This procedure provides a first-order depth estimate conditional on the supplied extent and boundary constraints. Propagation initialization, continuity refinement, and final depth calculation are sequential stages. Grid handling, screening thresholds, neighborhoods, and iteration settings are specified in Section~\ref{sec:eval_protocol}.

\clearpage
\section{Experiments and Results}

\subsection{Experimental configuration and evaluation protocol}\label{sec:eval_protocol}

\subsubsection{Datasets and data preparation}\label{sec:experiment_data}

The experimental data serve three purposes: learning bi-temporal flood changes, comparing flood-change detection across the three study areas, and evaluating depth within a supplied extent. The identified training snapshot combines S1GFloods \citep{saleh2024dam}, VarFloods \citep{du2026high}, and GF3 Henan samples, with membership summarized in Table~\ref{tab:training_split}. It contains 4,737 training pairs and 1,593 validation pairs, combining two public data sources with a small GF3 subset. GF3 Henan samples occur in both splits, and their spatial relationship to the Zhengzhou evaluation domain remains unresolved. The Zhengzhou comparison is therefore not defined as an independent cross-sensor transfer test.

\begin{table}[htbp]
\centering
\captionsetup{name=Table}
\caption{Training and validation dataset composition.}
\label{tab:training_split}
\small
\begin{tabular}{lrr}
\toprule
Data source & Training pairs & Validation pairs \\
\midrule
S1GFloods & 3,664 & 1,256 \\
VarFloods & 1,059 & 333 \\
GF3 Henan & 14 & 4 \\
\midrule
Total & 4,737 & 1,593 \\
\bottomrule
\end{tabular}
\end{table}


Learning samples use \(256\times256\) image tiles, a three-channel representation of the available SAR band, and binary labels for newly inundated areas. Normalization statistics come from the selected dataset. The semantic branch separately restores the image range and converts it to the normalization required by the pretrained model, as described in Section~\ref{sec:experiment_implementation}. Scene-specific preprocessing versions, label sources, invalid-pixel rules, and registration records still require association with the evaluated products; tile dimensions and input channels alone do not establish that association.

The scene comparison uses bi-temporal SAR imagery and flood-change reference labels for Zhengzhou, Zhuozhou, and Guangxi (Fig.~\ref{fig:gf3_eval_data}). Zhengzhou and Zhuozhou use GF3 imagery, while Guangxi uses LT-1B imagery. Observation parameters are listed in Table~\ref{tab:gf3_eval}, and terrain data are described in Section~\ref{sec:depth_data}. Zhengzhou and Zhuozhou use image pairs from different years, whereas the Guangxi pair was acquired in May and July 2026. The reference labels support spatial comparison with predictions, and the red boxes in the Zhengzhou and Zhuozhou panels indicate local inspection areas.

\begin{figure}[H]
  \centering
  \includegraphics[width=\textwidth,keepaspectratio]{figure/figure6.jpg}
  \caption{Bi-temporal SAR imagery and flood-change reference labels for the three study areas. Panels (a1)--(a3), (b1)--(b3), and (c1)--(c3) correspond to Zhengzhou, Zhuozhou, and Guangxi, respectively. Within each group, indices 1, 2, and 3 denote pre-event SAR imagery, post-event SAR imagery, and flood-change reference labels. Zhengzhou and Zhuozhou use GF3 data, while Guangxi uses LT-1B data. Red boxes indicate local inspection areas.}
  \label{fig:gf3_eval_data}
\end{figure}

Depth data follow a separate evaluation arrangement. The GF3 depth illustrations use prescribed inundation supports and the FABDEM data introduced in Section~\ref{sec:depth_data}. The Brazos river--floodplain comparison uses existing CFDepth and FwDET depth products and a reference raster attributed to USGS. It evaluates depth products within a supplied domain, independently of HarmoSSM flood detection. The reference product, vertical datum, exact support, and generating implementation remain incompletely linked; conclusions are therefore limited to the supplied products.

\subsubsection{Implementation and optimization settings}\label{sec:experiment_implementation}

Table~\ref{tab:method_config} consolidates the HarmoSSM representation and reconstruction settings, training selection, and CFDepth parameters. Model and depth parameters describe the current implementation, while optimization and selection values come from the identified training run. The association between that run and all GF3 comparison entries remains to be verified. The following input conversion, loss definitions, and inference settings are presented on this basis, without extending one training record to a common protocol for every comparison method.

\begin{table*}[htbp]
\centering
\captionsetup{name=Table}
\caption{Implementation settings for HarmoSSM and CFDepth.}
\label{tab:method_config}
\small
\setlength{\tabcolsep}{4pt}
\renewcommand{\arraystretch}{1.08}
\begin{tabularx}{\linewidth}{@{}p{0.23\linewidth}X@{}}
\toprule
Processing stage & Configuration \\
\midrule
\multicolumn{2}{@{}l}{\textit{HarmoSSM representation and reconstruction}} \\
Shared encoder & EfficientNet-B2; FPN: 128 channels at P1--P5 \\
Semantic branch & Frozen DINOv3 ViT-S/16; layers 5, 8, 11; input $512\times512$; fusion at P3--P5 \\
Paired harmonization & Calibration at P1/P2; local alignment at P1--P3; window 5; $\epsilon=10^{-5}$. P1: 2 heads, 4 samples, 2 offset groups; P2/P3: 4 heads, 9 samples, 4 groups \\
CQI & Five independent blocks; 128 channels, 16 queries and 4 heads per scale; FFN expansion 2 \\
OSCD & 128-channel fused context; 384-channel scan input; 4 scan directions; state dimension 1; FFN expansion 4 \\
\midrule
\multicolumn{2}{@{}l}{\textit{Training and selection}} \\
Input and batch & $256\times256$ tiles; batch size 12; bfloat16 mixed precision \\
Optimization & AdamW; base learning rate $10^{-4}$; head multiplier 2; weight decay $5\times10^{-4}$; cosine decay; 80 epochs \\
Checkpoint and threshold & Validation flood IoU; selected epoch 77 and threshold 0.80; search range 0.05--0.95 with step 0.01 \\
\midrule
\multicolumn{2}{@{}l}{\textit{CFDepth}} \\
Boundary screening & Maximum slope $5^\circ$; local-median radius 3 pixels; deviation multiplier 3; minimum tolerance 0.25 m \\
Anchors and initialization & Pairing radius 2 pixels; eight-connected propagation up to 10 km; terrain lower bound $Z+0.01$ m; low-confidence fallback where no anchor is reachable \\
Continuity refinement & Center-inclusive $3\times3$ window; relaxation 0.85; 50 iterations; fixed anchors and terrain lower bound \\
\bottomrule
\end{tabularx}
\end{table*}


The shared EfficientNet-B2 encoder and FPN produce five 128-channel levels, \(P1\)--\(P5\), at spatial reductions of 2, 4, 8, 16, and 32. Frozen DINOv3 ViT-S/16 uses a separate input conversion. For dataset-normalized input \(\mathbf{x}^t=(\mathbf{r}^t-\boldsymbol\mu_d)/\boldsymbol\sigma_d\), the image range is restored before applying the fixed ImageNet normalization associated with LVD pretraining:
\begin{equation}\label{eq:dino_input}
  \tilde{\mathbf{r}}^{t}
  =
  \operatorname{clip}\!\left(
  \mathbf{x}^{t}\odot\boldsymbol{\sigma}_{d}+\boldsymbol{\mu}_{d},0,1
  \right),
  \qquad
  \mathbf{x}_{\mathrm{DINO}}^{t}
  =
  \frac{\tilde{\mathbf{r}}^{t}-\boldsymbol{\mu}_{\mathrm{IM}}}
       {\boldsymbol{\sigma}_{\mathrm{IM}}}.
\end{equation}
The semantic input is then resized to \(512\times512\), distinct from the \(256\times256\) learning tile. Semantic features from layers 5, 8, and 11 are adapted to \(P3\), \(P4\), and \(P5\), respectively, and fused with CNN features through concatenation, convolution, and CBAM before HA; \(P1/P2\) retain CNN representations. The HA stabilizer, local-alignment heads and sampling points, and scale-specific CQI query settings are listed in Table~\ref{tab:method_config}.

OSCD organizes deep change information according to Eqs.~\ref{eq:oscd_deep_fusion} and~\ref{eq:oscd_scan_input}. Right and bottom replication padding makes the context grid divisible by four when necessary. The three \(\mathbf E_i\) receive the same padding, which is cropped after reconstructing \(\mathbf R_3\). Convolution \(\mathcal C_2\) uses kernel size 3 and stride 2, while \(\mathcal C_4\) uses kernel size 5 and stride 4; each branch is reduced to 128 channels. For an unpadded divisible \(H\times W\) input, \(\mathbf G\in\mathbb R^{384\times H/32\times W/32}\). The context module comprises positional encoding, four-direction selective state-space scanning, and a depthwise feed-forward update, followed by context aggregation and detail recovery according to Eqs.~\ref{eq:oscd_context_reconstruction} and~\ref{eq:oscd_reconstruction}.

Training constrains detection outputs through pixel classification and foreground overlap. For \(N\) labeled pixels and target-class probability \(p_{n,y_n}\), focal loss \citep{lin2017focal} uses class weights \(w_0=0.25,w_1=0.75\):
\begin{equation}\label{eq:focal_loss}
  \mathcal{L}_{\mathrm{focal}}
  =
  -\frac{1}{N}
  \sum_{n=1}^{N}
  w_{y_n}\left(1-p_{n,y_n}\right)^{\gamma}
  \log p_{n,y_n},
  \qquad \gamma=2.
\end{equation}
Tversky loss \citep{salehi2017tversky} computes soft TP, FP, and FN from foreground probabilities:
\begin{equation}\label{eq:tversky_loss}
 \mathcal{L}_{\mathrm{Tv}}=1-\frac{1}{B}\sum_{b=1}^{B}
 \frac{\mathrm{TP}_b+\epsilon_{\mathrm{Tv}}}{\mathrm{TP}_b+\alpha(e)\mathrm{FP}_b+\beta(e)\mathrm{FN}_b+\epsilon_{\mathrm{Tv}}},\qquad\alpha(e)=1-\beta(e).
\end{equation}
Here, \(B\) is the batch size, soft counts are computed per image, and \(\epsilon_{\mathrm{Tv}}=10^{-6}\). The false-negative weight \(\beta(e)\) decreases linearly from 0.70 at epoch 1 to 0.55 at epoch 20 and remains constant thereafter. The main head and five auxiliary heads share the same loss formulation:
\begin{equation}\label{eq:segmentation_loss}
  \mathcal{L}_{\mathrm{seg}}
  =
  \frac{1}{2}\mathcal{L}_{\mathrm{focal}}^{\mathrm{main}}
  +
  \mathcal{L}_{\mathrm{Tv}}^{\mathrm{main}}
  +
  \sum_{i=1}^{5}\omega_i(e)
  \left(
  \frac{1}{2}\mathcal{L}_{\mathrm{focal}}^{(i)}
  +
  \mathcal{L}_{\mathrm{Tv}}^{(i)}
  \right).
\end{equation}
The auxiliary ratios \((1,1,0.5,0.25,0.25)\) are normalized by their sum. Their common multiplier is 1 through epoch 5, decreases linearly to 0.5 at epoch 20, and remains 0.5 thereafter. Shallow support \(s=\max(p_{D1},p_{D2})\) is detached, and \(\mathcal L_{\mathrm{sup}}\) is the mean of \(\max(s-p_{\mathrm{main}},0)\) over \(s>0.60\). The coarse term averages \(\max(p_j-s-0.15,0)\) over \(s<0.35\), then over \(j\in\{\mathrm{main},D4,D5\}\); empty masks contribute zero. The complete objective is given in Eq.~\ref{eq:total_loss}. Both consistency weights are zero through epoch 6, increase linearly to 0.03 and 0.02 at epoch 15, and remain fixed. This corresponds to a five-epoch warm-up followed by a ten-epoch ramp whose first value is zero. These settings specify the loss implementation and do not establish an independently measured component benefit.

The implementation remains stored under the HA-CQI identifier. The associated record uses the 80-epoch optimization settings in Table~\ref{tab:method_config}, random seed 1, and shared bi-temporal radiometric augmentation. Validation flood IoU jointly selects the checkpoint and threshold, with epoch 77 and 0.80 selected in the record. This validation selection is distinct from GF3 scene-comparison results. The eight comparison methods are FC-Siam-Diff \citep{daudt2018fully}, HANet \citep{han2023hanet}, SNUNet \citep{fang2022snunet}, BIT \citep{chen2021remote}, ChangeFormer \citep{bandara2022changeformer}, CGNet \citep{han2023change}, TTP \citep{chen2023ttp}, and ChangeDINO \citep{cheng2025changedino}. Together with HarmoSSM, they form the nine-method set in Table~\ref{tab:main_comparison}. That table transcribes author-supplied results; row-specific checkpoints, predictions, thresholds, and valid domains have not been independently verified. The configuration table therefore does not establish identical training settings across all methods.

CFDepth maps the extent to the DEM grid by maximum-value aggregation and intersects it with the valid DEM mask. Candidate exterior pixels are eight-neighbor adjacent and screened using the slope and local-deviation thresholds in Table~\ref{tab:method_config}. Local deviations are constructed sequentially from candidate elevation medians, pointwise absolute deviations, and local medians of the deviation field. Radius-two interior and exterior boundary neighborhoods provide paired anchors. Nearest reachable anchors propagate within eight-connected support up to 10 km; pixels beyond that distance or without reachable anchors fall back to \(Z+0.01\) m and retain low-confidence flags. Continuity refinement uses a center-inclusive \(3\times3\) window, applies the terrain lower bound, and restores fixed anchors at each update. Locations outside the support or without valid DEM data are NoData. CFDepth does not enter the detection model's loss optimization.

\subsubsection{Evaluation metrics and protocols}\label{sec:experiment_metrics}

Flood-change mapping treats flood as the positive class, with \(TP\), \(FP\), \(FN\), and \(TN\) denoting confusion counts over the valid labeled domain. Precision (P) measures the correct fraction of predicted flood, recall (R) measures the detected fraction of reference flood, F1 and IoU summarize detection performance and foreground overlap, and OA measures correctness across all valid pixels:
\begin{align}
 \mathrm{Precision}&=\frac{TP}{TP+FP}, & \mathrm{Recall}&=\frac{TP}{TP+FN},\label{eq:precision}\\
 \mathrm{F1}&=\frac{2TP}{2TP+FP+FN}, & \mathrm{IoU}&=\frac{TP}{TP+FP+FN},\label{eq:f1}\\
 \mathrm{OA}&=\frac{TP+TN}{TP+FP+FN+TN}.\label{eq:oa}
\end{align}
All five metrics are reported as percentages. Because background prevalence affects OA, flood-region overlap is assessed primarily through IoU and F1, with P and R explaining commission--omission trade-offs. This chapter uses only the author-designated result sheet dated 2 September 2026. Numerical transcription and relationships among metrics can be checked, but those checks do not replace verification of the original predictions, labels, thresholds, and valid evaluation domains.

Depth coverage and error use different statistical domains. In the existing Brazos evaluation, reference resampling uses a valid-area fraction threshold of 0.50, and finite positive-depth pixels define the reference domain:
\begin{equation}\label{eq:valid_depth_ratio}
 \Omega_r=\{p:\ d_r(p)>0.01\,\mathrm{m},\ d_r(p)\text{ finite}\},\qquad
 R_j=100\frac{|V_j|}{|\Omega_r|},
\end{equation}
Here, \(V_j=\{p:\ p\in\Omega_r,\ d_j(p)>0.01\,\mathrm{m},\ d_j(p)\text{ finite}\}\) is the valid set for method \(j\). The ratio \(R_j\) measures positive-depth coverage within this reference domain. MAE, RMSE, Bias, squared Pearson correlation \(r^2\), and Nash--Sutcliffe efficiency (NSE) are computed on \(V_{\mathrm{CFDepth}}\cap V_{\mathrm{FwDET}}\). NoData is not replaced by zero, avoiding the treatment of unavailable estimates as zero-depth errors.

Bias is prediction minus reference. Pearson \(r^2\) measures linear association, whereas NSE compares squared error with reference variation about its mean. Lower aggregate error, broader coverage, and closer spatial agreement must therefore be assessed separately. The evaluation script does not itself remove permanent water; any earlier masking remains part of the unresolved product lineage. This protocol keeps depth-product comparison separate from flood-extent detection.

\subsection{Quantitative flood-change mapping results}

To compare bi-temporal flood-extent delineation, Table~\ref{tab:main_comparison} summarizes P, R, F1, IoU, and OA for nine methods in Zhengzhou and Zhuozhou. Values are transcribed from the result sheet dated 2 September 2026, supplied by the authors and identified by them as measured results; the original predictions have not been reevaluated. The following comparison uses IoU and F1 as the main overlap measures and P and R to examine commission--omission trade-offs. All numerical comparisons are confined to the methods and results in this table.

\begin{table*}[htbp]
  \centering
  \caption{Quantitative comparison of urban flood extent delineation on the Zhengzhou and Zhuozhou evaluation scenes.}
  \label{tab:main_comparison}
  \scriptsize
  \setlength{\tabcolsep}{4pt}
  \begin{adjustbox}{width=\textwidth}
    \begin{tabular}{lcccccccccc}
      \toprule
      \multirow{2}{*}{Methodology} & \multicolumn{5}{c}{Zhengzhou} & \multicolumn{5}{c}{Zhuozhou} \\
      \cmidrule(lr){2-6}\cmidrule(lr){7-11}
      & IoU$\uparrow$ & F1$\uparrow$ & Precision$\uparrow$ & Recall$\uparrow$ & OA$\uparrow$
      & IoU$\uparrow$ & F1$\uparrow$ & Precision$\uparrow$ & Recall$\uparrow$ & OA$\uparrow$ \\
      \midrule
      FC-Siam-Diff & 62.59 & 76.99 & 80.05 & 74.16 & 99.41 & 70.87 & 82.95 & 86.95 & 79.30 & 97.11 \\
      HANet & 49.10 & 65.86 & 61.86 & 70.41 & 99.03 & 62.95 & 77.26 & 81.26 & 73.63 & 96.15 \\
      SNUNet & 50.39 & 67.01 & 71.01 & 63.44 & 99.17 & 70.09 & 82.41 & 86.41 & 78.77 & 97.02 \\
      BIT & 54.72 & 70.73 & 74.73 & 67.14 & 99.26 & 72.23 & 83.88 & 87.88 & 80.23 & 97.26 \\
      ChangeFormer & 50.99 & 67.54 & 71.54 & 63.97 & 99.18 & 67.21 & 80.39 & 84.39 & 76.75 & 96.68 \\
      CGNet & 62.76 & 77.12 & 73.12 & 81.58 & 99.35 & 82.20 & 90.23 & 94.23 & 86.56 & 98.34 \\
      TTP & 54.44 & 70.50 & 74.50 & 66.91 & 99.25 & 79.59 & 88.63 & 92.63 & 84.96 & 98.07 \\
      ChangeDINO & 33.65 & 50.35 & 54.17 & 47.03 & 98.76 & 71.72 & 83.53 & 87.53 & 79.88 & 97.20 \\
      \textbf{Ours (HA-CQI)} & \textbf{73.38} & \textbf{84.65} & \textbf{80.65} & \textbf{89.06} & \textbf{99.57} & \textbf{86.23} & \textbf{92.61} & \textbf{96.61} & \textbf{88.92} & \textbf{98.74} \\
      \bottomrule
    \end{tabular}
  \end{adjustbox}
\end{table*}


The main pattern in the Zhengzhou entries is the combination of high flood recall and higher overall overlap. HarmoSSM has an IoU of 74.23\% and an F1 of 85.21\%, exceeding CGNet's 67.50\% and 80.60\% by 6.73 and 4.61 percentage points, respectively. Its recall of 91.59\% is close to HANet's 92.74\%, while its precision of 79.66\% is higher than HANet's 58.08\%. Conversely, ChangeFormer has a precision of 97.95\% but a recall of 49.24\%. These reported values show that a high correct fraction of predicted flood or a high coverage of reference flood alone does not ensure high overlap. In this comparison, HarmoSSM retains much of the reference flood while limiting commission relative to the high-recall comparator, yielding higher F1 and IoU.

The Zhuozhou entries instead emphasize differences in coverage under high precision. HarmoSSM has an IoU of 86.71\% and an F1 of 92.88\%, exceeding CGNet by 4.31 and 2.53 percentage points, respectively; its precision is 97.58\% and recall is 88.61\%. ChangeFormer has higher precision, 99.39\%, but recall of 75.87\%, while CGNet has precision and recall of 93.08\% and 87.78\%. HarmoSSM's overlap in this result sheet therefore corresponds to a combination of a high correct fraction of predictions and broader reference-flood coverage, rather than the highest precision alone. Its OA of 98.84\% is consistent with the foreground metrics, although OA cannot substitute for flood-region evaluation.

The supplied results place HarmoSSM highest in IoU and F1 in both scenes, while its relative precision and recall vary by scene. Zhengzhou emphasizes the balance between recall and commission, whereas Zhuozhou emphasizes coverage under high precision. These differences relate to the need to organize regional change information together with local positioning cues, providing quantitative context for the spatial illustrations that follow. The comparison does not separate the contributions of HCE, OSCD, or their internal operations, and it establishes neither statistical significance nor independent cross-sensor generalization.

\subsection{Qualitative comparison in complex scenes}

Qualitative comparison examines spatial errors that aggregate metrics cannot locate. Figures~\ref{fig:qual_zhengzhou} and~\ref{fig:qual_zhuozhou} present pre-event and post-event images, references, and archived outputs for Zhengzhou and Zhuozhou, respectively. Relative to the reference, white denotes TP, blue FN, red FP, and black TN. Output columns have not been linked to specific methods and checkpoints. The following interpretation therefore uses identifiable scene structures and error patterns without assigning the main-table method order to panel columns.

The Zhengzhou locations emphasize retention of small patches and narrow inundated strips within built-up and linear-feature backgrounds. The upper reference rows contain isolated patches, and the lower rows include narrow structures represented by relatively few pixels. Output panels show blue omissions and red commissions at different positions around these features. Omitting a few pixels can weaken the representation of a local inundated area, while commissions in adjacent background enlarge the predicted extent, making recall and precise localization relevant together. This observation explains the joint use of P, R, and overlap measures in Section~4.2, but the selected panels do not establish a detectable-width threshold or weak-contrast accuracy.

\begin{figure*}[!t]
  \centering
  \includegraphics[width=\textwidth,height=.76\textheight,keepaspectratio]{figure/figure7.jpg}
  \caption{Flood-change detection results in Zhengzhou: (a) pre-event SAR; (b) post-event SAR; (c) reference mask; (d)--(l) predictions. White, blue, red, and black denote TP, FN, FP, and TN, respectively.}
  \label{fig:qual_zhengzhou}
\end{figure*}

The Zhuozhou locations require both coverage of contiguous areas and retention of local separations. The second reference row contains a broad inundated region; subsequent rows show field margins, fragmented areas near settlements, and finer separations. Differences in interior coverage and boundary errors across outputs indicate that regional judgments must be translated into specific spatial extents: filling interior omissions and retaining real separations impose different requirements. Together with Zhengzhou, these examples place regional context and local position cues within the same reconstruction task. They explain scene-level mapping needs but do not constitute scene-wide morphology-stratified statistics or causal validation of model components.

\begin{figure*}[!t]
  \centering
  \includegraphics[width=\textwidth,height=.76\textheight,keepaspectratio]{figure/figure8.jpg}
  \caption{Flood-change detection results in Zhuozhou: (a) pre-event SAR; (b) post-event SAR; (c) reference mask; (d)--(l) predictions. White, blue, red, and black denote TP, FN, FP, and TN, respectively.}
  \label{fig:qual_zhuozhou}
\end{figure*}

\subsection{Conditional depth-estimation results}

Depth results evaluate conditional estimates within a supplied extent through spatial availability and error within a reference domain. The GF3 illustrations show the spatial distributions of existing depth products, whereas the Brazos comparison uses an independent reference raster and common-valid pixels for quantitative assessment. These data address different questions; neither provides joint accuracy validation of the full HarmoSSM detection and CFDepth estimation chain.

Figures~\ref{fig:depth_zhengzhou} and~\ref{fig:depth_zhuozhou} show GF3 products labeled FwDET and CFDepth, with differences in colored coverage and unresolved areas. Availability of a depth estimate depends on the supplied support, valid DEM, and usable boundary constraints, so these differences first indicate where the products provide information. Support masks and generating versions remain incompletely linked; no coverage percentage, smoothness metric, or photo-based accuracy is derived from the illustrations. Figure~\ref{fig:wse_gradient_cdf} retains the archived WSE-gradient diagnostic for layout review; its source statistics are unverified and are not used quantitatively.

\begin{figure*}[!t]
  \centering
  \includegraphics[width=\textwidth,height=.69\textheight,keepaspectratio]{figure/figure9.jpg}
  \caption{Flood-depth estimates in Zhengzhou: (a) FwDET; (b) CFDepth. Insets p1--p4 correspond to the photograph locations in Fig.~\ref{fig:dataset}; white denotes unresolved depth.}
  \label{fig:depth_zhengzhou}
\end{figure*}

\begin{figure*}[!t]
  \centering
  \includegraphics[width=\textwidth,height=.69\textheight,keepaspectratio]{figure/figure10.jpg}
  \caption{Flood-depth estimates in Zhuozhou: (a) FwDET; (b) CFDepth. Insets p5--p8 correspond to the photograph locations in Fig.~\ref{fig:dataset}; white denotes unresolved depth.}
  \label{fig:depth_zhuozhou}
\end{figure*}

\begin{figure*}[!t]
  \centering
  \includegraphics[width=\linewidth,height=0.70\textheight,keepaspectratio]{figure/figure11.jpg}
  \caption{Cumulative distributions of water-surface elevation (WSE) gradients for FwDET and CFDepth. Dashed lines indicate the corresponding medians.}
  \label{fig:wse_gradient_cdf}
\end{figure*}

The Brazos product comparison shows differences in coverage and common-domain error in the same direction (Table~\ref{tab:cfdepth_validation}; Fig.~\ref{fig:usgs_depth_validation}). The CFDepth-labeled product has positive-depth coverage of 95.810\%, compared with 41.910\% for FwDET. On common-valid pixels, its MAE and RMSE are 1.016 and 1.416 m, compared with 1.673 and 2.454 m, respectively. Bias is positive for both products, at 0.316 and 1.030 m. Thus, the CFDepth-labeled product supplies broader positive-depth coverage in this reference domain and has lower absolute error, squared error, and positive bias where both products provide valid depths. The additional coverage itself does not enter the common-domain error calculation.

\begin{table*}[htbp]
\centering
\captionsetup{name=Table}
\caption{Comparison of flood-depth estimates in the Brazos River floodplain.}
\label{tab:cfdepth_validation}
\small
\setlength{\tabcolsep}{5pt}
\begin{tabular}{lrrrrrr}
\toprule
Product label & Valid depth (\%) & MAE (m) & RMSE (m) & Bias (m) & Pearson $r^2$ & NSE \\
\midrule
FwDET & 41.910 & 1.673 & 2.454 & 1.030 & 0.053 & -7.493 \\
CFDepth & 95.810 & 1.016 & 1.416 & 0.316 & 0.150 & -1.828 \\
\bottomrule
\end{tabular}
\end{table*}


Lower aggregate error does not coincide with close spatial agreement in depth. Pearson \(r^2\) is 0.150 for CFDepth and 0.053 for FwDET, while NSE is \(-1.828\) and \(-7.493\), respectively. Low correlation and negative NSE indicate limited representation of reference depth variation. In particular, negative NSE means that squared error exceeds that of a constant prediction equal to the common-domain reference mean. The comparison therefore supports relative differences in coverage and error between the supplied products while limiting their interpretation as detailed depth distributions. The evaluated rasters have not been fully linked to the current implementation, so these differences cannot be attributed to specific CFDepth operations or generalized to measured urban flood-depth accuracy.

\begin{figure*}[!t]
  \centering
  \includegraphics[width=\textwidth,height=.43\textheight,keepaspectratio]{figure/figure12.jpg}
  \caption{Spatial distributions of flood depth in the Brazos River floodplain: (a) reference depth; (b) CFDepth; (c) FwDET.}
  \label{fig:usgs_depth_validation}
\end{figure*}

\section{Discussion}

\subsection{Landscape-dependent mapping and interpretation}

The two scenes expose different spatial requirements of flood mapping. In the Zhengzhou panels, a small number of pixels may represent an inundated patch beside a built structure or linear feature. In Zhuozhou, regional coverage must be maintained while resolving field margins and fragmented transitions. A single scene can therefore require both contextual interpretation and detailed delineation. These observations motivate the joint treatment of the two spatial scales, but selected windows cannot determine how frequently each difficulty occurs across a scene.

HarmoSSM implements this treatment by coordinating the paired feature basis before change interaction and retaining shallow detail during dense reconstruction. The directed difference preserves the sign of feature change, while the two state features provide context for its interpretation. These are properties of the design, not direct measurements of reduced false change or improved narrow-feature recall. AWCA-Net and RAFNet already address heterogeneous flood responses through different processing strategies \citep{du2026high,li2026rafnet}. The relevant distinction is how information is organized for the present binary change target; component-level benefit cannot be inferred from architectural differences or visual examples alone.

\subsection{Conditional value of depth information}

Depth adds information inside an inundation support, but its interpretation depends on the available constraints. Boundary-derived WSE can be useful where valid wet--dry interfaces and suitable terrain data are available; permanent-water interfaces and unobserved margins weaken those constraints. FwDET, FLEXTH, and RS-FloodXDepth address related boundary and extent limitations \citep{cohen2019floodwater,betterle2024water,tian2026floodxdepth}. CFDepth adopts screening, constrained initialization, and local continuity refinement within a supplied support. This formulation produces a terrain-consistent first-order surface, without enforcing water balance or observed hydraulic connections.

The Brazos product comparison illustrates why output coverage and depth accuracy require separate assessment. More positive-depth pixels do not necessarily imply more accurate depths, and lower common-valid error does not establish accurate spatial patterning. Negative NSE means that, on the evaluated common-valid set, squared error exceeds that of predicting its reference mean. This limits the quantitative interpretation even where aggregate error is lower than that of the compared product. The unresolved generating-version lineage further prevents attributing these differences to a specific CFDepth operation.

\subsection{Applicability and limitations}

SAR observability remains a physical limitation. Building geometry, layover, shadow, vegetation, roughness, and acquisition timing can obscure flooding or resemble non-flood change \citep{wagner2026gfm,tsyganskaya2018vegetation}. Feature calibration cannot guarantee recovery of an unobserved water signal. Likewise, label resolution and uncertain event timing affect which regions are counted as newly inundated \citep{zhang2025floodplanet}. The unresolved training/evaluation overlap and prediction provenance preclude an independent cross-sensor conclusion; dataset-dependent generalization requires evidence beyond a change of sensor name \citep{portales2025flood}.

CFDepth also inherits the spatial and vertical limits of the DEM. A narrow SAR detection may be represented by only part of a terrain cell, and retaining its support through maximum aggregation does not recover subcell topography. Erroneous supports change both the estimated domain and the available boundary elevations; this is a dependency of the formulation, not a measured perturbation result. A processing-chain study of Garonne floods also shows the influence of input extents and parameter choices on depth estimates \citep{travert2026uncertainty}, although its sensitivity magnitudes cannot be transferred to the present cases. Bridges, culverts, drainage, and exchange between image-disconnected regions are not represented. Pixels lacking reachable anchors use a terrain-relative fallback and should be distinguished from anchored estimates. Together with the unverified photograph timing and reference-product lineage, these limitations restrict the depth results to conditional information rather than surveyed flood depths or event dynamics.

\section{Conclusion}

HarmoSSM addresses newly inundated land in complex landscapes by combining comparable bi-temporal SAR representations, state and change evidence, and regional context with local reconstruction. CFDepth extends a supplied inundation extent to a first-order depth field using screened terrain constraints. The formulation connects the location of new inundation with conditional information inside that extent while preserving their different observation and evaluation requirements.

The Zhengzhou and Zhuozhou panels illustrate the coexistence of localized, narrow, fragmented, and extensive inundation. They support the relevance of the mapping problem but do not establish morphology-specific improvements or component effects. In the separate Brazos product comparison, broader positive-depth coverage and lower aggregate error coexist with low correlation and negative NSE. These findings constrain interpretation to the displayed spatial cases and supplied depth products; they do not establish independent sensor transfer, surveyed urban depth accuracy, or hydrodynamic fidelity.

\section*{Acknowledgements}

\begin{sloppypar}
This work was supported by the National Key R\&D Program of China (Grant No. 2024YFC3808602). The authors thank Prof. Shuliang Zhang from Nanjing Normal University for his valuable assistance with this research.
\end{sloppypar}

\bibliographystyle{elsarticle-harv}
\bibliography{reference}

\end{document}
